\chapter{Testy}
\label{chap:testy}

\section{Zakres testow}
\label{sec:zakres-testow}

Testy jednostkowe obejmuja parser, tworzenie stempli, obsluge bledow oraz solver. Osobne testy sprawdzaja poprawne dane wejsciowe i przypadki, w ktorych program powinien zglosic blad.

\section{Testy parsera}
\label{sec:testy-parsera}

Parser jest testowany na pojedynczych elementach oraz na pliku zawierajacym wszystkie obslugiwane typy elementow. Pozwala to sprawdzic, czy program poprawnie rozpoznaje elementy, parametry oraz komende analizy.

\section{Testy solvera}
\label{sec:testy-solvera}

Testy solvera porownuja wynik rozwiazania ukladu z oczekiwanymi wartosciami dla prostych obwodow. Takie przypadki testowe sa latwe do zweryfikowania recznie i dobrze wykrywaja bledy w skladaniu macierzy.

\section{Testy bledow stempli}
\label{sec:testy-bledow-stempli}

Bledy parametrow stempli sa sprawdzane oddzielnie od testow parsera i solvera. Dzieki temu testy pozostaja czytelne i kazdy zestaw sprawdza jedna warstwe programu.
